\documentclass[11pt]{article}
% ===========================================================================
%  Shared preamble for the AP Calculus AB exam-preparation worksheets.
%  Input right after \documentclass. Files beginning with "_" are not compiled
%  on their own. Navy and gold, to match the course's lesson pages.
% ===========================================================================
\usepackage[letterpaper,top=2.3cm,bottom=2.2cm,left=1.7cm,right=1.7cm,headheight=28pt,footskip=30pt]{geometry}
\usepackage{amsmath,amssymb}
\usepackage[T1]{fontenc}
\usepackage[utf8]{inputenc}
\usepackage{lmodern}
\usepackage{xcolor}
\usepackage[most]{tcolorbox}
\usepackage{enumitem}
\usepackage{fancyhdr}
\usepackage{array}
\usepackage{tikz}
\usepackage{pgfplots}
\pgfplotsset{compat=1.18}

\definecolor{navy}{HTML}{161E54}
\definecolor{ink}{HTML}{27357E}
\definecolor{gold}{HTML}{B98A00}
\definecolor{goldbg}{HTML}{FFFDF3}
\definecolor{inkbg}{HTML}{F6F7FD}
\definecolor{rule}{HTML}{C9CFF0}
\definecolor{okgreen}{HTML}{1B7A44}
\definecolor{okbg}{HTML}{F1F8F3}

\pagestyle{fancy}
\fancyhf{}
\fancyhead[L]{\footnotesize Name: \underline{\hspace{4cm}}}
\fancyhead[C]{\textbf{\color{navy}Integration Academy}\\[-2pt]{\scriptsize Calculus AB \textperiodcentered{} AP\textsuperscript{\textregistered} Exam Preparation}}
\fancyfoot[L]{\tiny\color{black!55}\parbox[t]{0.78\textwidth}{AP\textsuperscript{\textregistered} is a registered trademark of the College Board, which was not involved in the production of, and does not endorse, this material. \textcopyright{} Integration Academy.}}
\fancyfoot[R]{\scriptsize\thepage}
\renewcommand{\headrulewidth}{1.2pt}
\renewcommand{\headrule}{\hbox to\headwidth{\color{gold}\leaders\hrule height \headrulewidth\hfill}}
\renewcommand{\footrulewidth}{0pt}
\setlength{\parindent}{0pt}
\setlength{\parskip}{4pt}

% Title block: \wstitle{code}{title}{unit line}
\newcommand{\wstitle}[3]{\begin{center}{\small\bfseries\color{gold}WORKSHEET #1}\\[3pt]{\LARGE\bfseries\color{navy}#2}\\[3pt]{\small\color{black!60}#3}\end{center}\vspace{2pt}}
% Section heading
\newcommand{\sech}[1]{\par\vspace{10pt}{\large\bfseries\color{navy}#1}\par\nobreak\vspace{2pt}{\color{rule}\hrule height0.8pt}\vspace{7pt}}
% A level tag printed beside a group of questions
\newcommand{\level}[1]{\par\vspace{6pt}{\footnotesize\bfseries\color{gold}\MakeUppercase{#1}}\par\nobreak\vspace{2pt}}

\newtcolorbox{ideabox}{colframe=ink,colback=inkbg,coltitle=white,colbacktitle=ink,fonttitle=\bfseries\small,title={Key ideas},arc=2pt,boxrule=0.8pt,left=8pt,right=8pt}
\newtcolorbox{example}[1]{colframe=rule,colback=white,coltitle=navy,colbacktitle=inkbg,fonttitle=\bfseries\small,title={#1},arc=2pt,boxrule=0.8pt,borderline west={3pt}{0pt}{ink},left=10pt,right=8pt}
\newtcolorbox{frbox}[1]{colframe=gold,colback=goldbg,coltitle=white,colbacktitle=gold,fonttitle=\bfseries\small,title={#1},arc=2pt,boxrule=0.8pt,breakable,left=8pt,right=8pt}
\newtcolorbox{solbox}{colframe=okgreen,colback=okbg,coltitle=white,colbacktitle=okgreen,fonttitle=\bfseries\small,title={Worked solutions},arc=2pt,boxrule=0.8pt,breakable,left=8pt,right=8pt}
\newtcolorbox{scorebox}[1]{colframe=navy,colback=white,coltitle=white,colbacktitle=navy,fonttitle=\bfseries\small,title={#1},arc=2pt,boxrule=0.8pt,breakable,left=8pt,right=8pt}

\newcommand{\soln}{\par\smallskip\textbf{\color{ink}Solution.}\ }
\newcommand{\result}[1]{\par\smallskip\textbf{#1}}
% blank working space of a given height
\newcommand{\work}[1]{\par\vspace{#1}}
% a point value in the margin of a scoring guide
\newcommand{\pts}[1]{\hfill{\footnotesize\bfseries\color{gold}[#1]}}

% Coordinate grid for graph questions: \begin{apgraph}{xmin}{xmax}{ymin}{ymax} ... \end{apgraph}
\newenvironment{apgraph}[4]{\begin{center}\begin{tikzpicture}\begin{axis}[
  axis lines=middle, xmin=#1, xmax=#2, ymin=#3, ymax=#4,
  grid=both, grid style={black!12}, xtick distance=1, ytick distance=1,
  ticklabel style={font=\scriptsize}, xlabel={$x$}, ylabel={$y$},
  width=10cm, height=7.2cm, samples=80, clip=false,
  every axis plot/.append style={line width=1.1pt, color=ink}]}{\end{axis}\end{tikzpicture}\end{center}}
\newcommand{\closeddot}[2]{\addplot[only marks, mark=*, mark size=2.2pt, color=ink] coordinates {(#1,#2)};}
\newcommand{\opendot}[2]{\addplot[only marks, mark=*, mark size=2.2pt, color=ink, mark options={fill=white}] coordinates {(#1,#2)};}

% Multiple-choice options: four across, or stacked when the options are long.
\newcommand{\opts}[4]{\par\vspace{3pt}\begin{tabular}{@{}p{0.235\linewidth}p{0.235\linewidth}p{0.235\linewidth}p{0.235\linewidth}@{}}(A)\ #1 & (B)\ #2 & (C)\ #3 & (D)\ #4\end{tabular}\par\vspace{9pt}}
\newcommand{\optsv}[4]{\par\vspace{3pt}\begin{tabular}{@{}l@{\ \ }p{0.9\linewidth}@{}}(A) & #1\\(B) & #2\\(C) & #3\\(D) & #4\end{tabular}\par\vspace{9pt}}

\fancyhead[R]{\footnotesize\color{black!60}Worksheet 1.10}
\begin{document}

\wstitle{1.10}{Guaranteed Values: The Intermediate Value Theorem}{Unit 1: Limits and Continuity}

\begin{ideabox}
\begin{itemize}[leftmargin=*,itemsep=2pt,topsep=2pt]
  \item \textbf{IVT.} If $f$ is continuous on $[a,b]$ and $N$ is between $f(a)$ and $f(b)$, then $f(c)=N$ for at least one $c$ in $(a,b)$.
  \item A justification has three parts: continuity, the inequality with numbers, and the theorem named.
  \item The theorem guarantees \emph{at least one} value; it does not locate it or count them.
  \item To show $f(c)=g(c)$ for some $c$, apply the theorem to $h=f-g$.
\end{itemize}
\end{ideabox}

\sech{Worked examples}
\begin{example}{Example 1 \textperiodcentered{} A root in an interval}
Show that $x^3-4x+1=0$ has a solution in $(0,1)$.\soln Let $f(x)=x^3-4x+1$, a polynomial, so $f$ is continuous on $[0,1]$. $f(0)=1$ and $f(1)=-2$, so $f(1)<0<f(0)$.\result{By the Intermediate Value Theorem there is a $c$ in $(0,1)$ with $f(c)=0$.}
\end{example}

\begin{example}{Example 2 \textperiodcentered{} Counting from a table}
A function $f$ is continuous on $[0,9]$.\begin{center}\renewcommand{\arraystretch}{1.25}\begin{tabular}{|c|c|c|c|c|c|}\hline $x$ & $0$ & $2$ & $4$ & $7$ & $9$ \\ \hline $f(x)$ & $5$ & $-1$ & $2$ & $6$ & $-3$ \\ \hline\end{tabular}\end{center}What is the least number of zeros of $f$ on $[0,9]$?\soln The sign changes on $[0,2]$, $[2,4]$ and $[7,9]$, but not on $[4,7]$.\result{At least $3$.}
\end{example}

\begin{example}{Example 3 \textperiodcentered{} Two functions}
$f$ and $g$ are continuous, with $f(2)=1$, $g(2)=6$, $f(5)=9$, $g(5)=4$. Show that $f(c)=g(c)$ for some $c$ in $(2,5)$.\soln Let $h=f-g$, continuous on $[2,5]$. $h(2)=-5<0$ and $h(5)=5>0$.\result{By the Intermediate Value Theorem $h(c)=0$ for some $c$ in $(2,5)$, so $f(c)=g(c)$.}
\end{example}

\newpage
\sech{Practice}

\level{Core}
\begin{enumerate}[leftmargin=*,itemsep=2pt]
  \item State the hypothesis of the Intermediate Value Theorem.\work{2.3cm}
  \item $f$ is continuous on $[0,3]$ with $f(0)=-2$ and $f(3)=7$. Which of the values $0$, $5$, $8$ must $f$ take on $(0,3)$?\work{2.3cm}
  \item Show that $x^5+x-1=0$ has a solution in $(0,1)$.\work{2.3cm}
  \item Explain why the Intermediate Value Theorem cannot be applied to $f(x)=\dfrac{1}{x-1}$ on $[0,2]$.\work{2.3cm}
\end{enumerate}

\level{Exam level}
\begin{enumerate}[leftmargin=*,itemsep=2pt,start=5]
  \item Show that the equation $e^x=4-x^2$ has a solution in $(1,2)$.\work{2.8cm}
  \item A differentiable function $P$ has $P(1)=12$ and $P(4)=-3$. Must there be a $t$ in $(1,4)$ with $P(t)=5$? Justify.\work{2.8cm}
  \item $f$ is continuous on $[1,6]$ with $f(1)=4$ and $f(6)=4$. Must $f$ have a zero in $(1,6)$? Explain.\work{2.8cm}
  \item For which values of $k$ does the theorem guarantee that $f(x)=x^2+x+k$ has a zero in $(0,1)$?\work{2.8cm}
\end{enumerate}

\begin{frbox}{Question 9 \textperiodcentered{} Free-response practice}
The temperature in a greenhouse, $T(t)$ $^\circ$C, is a continuous function of time $t$ hours.\begin{center}\renewcommand{\arraystretch}{1.25}\begin{tabular}{|c|c|c|c|c|}\hline $t$ & $0$ & $5$ & $10$ & $20$ \\ \hline $T(t)$ & $18$ & $25$ & $21$ & $30$ \\ \hline\end{tabular}\end{center}
\begin{enumerate}[label=(\alph*),leftmargin=*,itemsep=2pt]
  \item Must there be a time $t$, $0<t<5$, at which $T(t)=23$? Justify your answer.\work{2.6cm}
  \item What is the least number of times in $[0,20]$ at which $T(t)=23$? Show your reasoning.\work{2.6cm}
  \item A student claims the table shows the temperature never fell below 18 $^\circ$C. Is the claim justified?\work{2.6cm}
\end{enumerate}
\end{frbox}

\newpage
\begin{solbox}
\begin{enumerate}[leftmargin=*,itemsep=4pt]
  \item $f$ is continuous on the closed interval $[a,b]$.
  \item $0$ and $5$ (both lie between $-2$ and $7$). Not necessarily $8$.
  \item $f(x)=x^5+x-1$ is continuous on $[0,1]$; $f(0)=-1<0<1=f(1)$. By the IVT there is a $c$ in $(0,1)$ with $f(c)=0$.
  \item $f$ is not continuous on $[0,2]$: it is undefined at $x=1$.
  \item Let $h(x)=e^x+x^2-4$, continuous on $[1,2]$. $h(1)=e-3<0$ and $h(2)=e^2>0$. By the IVT $h(c)=0$ for some $c$ in $(1,2)$.
  \item Yes. $P$ is differentiable, hence continuous, on $[1,4]$, and $P(4)=-3<5<12=P(1)$. By the IVT such a $t$ exists.
  \item Not necessarily. $0$ is not between $f(1)$ and $f(6)$, so the theorem guarantees nothing.
  \item $f(0)=k$ and $f(1)=2+k$. Zero lies between them when $k<0<2+k$: $-2<k<0$.
  \item \begin{enumerate}[label=(\alph*),leftmargin=*,itemsep=3pt]
    \item Yes. $T$ is continuous on $[0,5]$ and $T(0)=18<23<25=T(5)$. By the Intermediate Value Theorem there is a time $t$ in $(0,5)$ with $T(t)=23$.
    \item 23 lies between the endpoint values on $[0,5]$ (18 and 25), on $[5,10]$ (25 and 21) and on $[10,20]$ (21 and 30). At least $3$ times.
    \item No. The table gives the temperature at four times only; between them $T$ may have dipped below 18 and returned.
  \end{enumerate}
\end{enumerate}
\end{solbox}

\end{document}
